\documentclass[12pt, a4paper]{article}
\usepackage[utf8]{inputenc}
\usepackage[T1]{fontenc}
\usepackage[ngerman]{babel}
\usepackage{amsmath, amssymb, amsthm}
\usepackage{geometry}
\usepackage{graphicx}
\usepackage{xcolor}
\usepackage{tcolorbox}
\usepackage{enumitem}
\usepackage{listings}
\usepackage{fancyhdr}
\usepackage{titlesec}
\usepackage{mdframed}
\usepackage{tikz}
\usepackage{pgfplots}
\pgfplotsset{compat=1.18}

\geometry{a4paper, margin=2.5cm}

\definecolor{theoryblue}{RGB}{30, 90, 160}
\definecolor{exerciseorange}{RGB}{200, 80, 0}
\definecolor{solutiongreen}{RGB}{0, 120, 60}
\definecolor{hintgray}{RGB}{100, 100, 100}
\definecolor{codebg}{RGB}{245, 245, 245}

% Theory box
\tcbuselibrary{skins, breakable}
\newtcolorbox{theorybox}[1]{
  colback=theoryblue!5,
  colframe=theoryblue,
  fonttitle=\bfseries\color{white},
  title=#1,
  breakable
}

% Exercise box
\newtcolorbox{exercisebox}[1]{
  colback=exerciseorange!5,
  colframe=exerciseorange,
  fonttitle=\bfseries\color{white},
  title=Aufgabe #1,
  breakable
}

% Solution space box
\newtcolorbox{solutionspace}[1][6cm]{
  colback=white,
  colframe=black!20,
  fonttitle=\itshape\color{hintgray},
  title=Rechenraum,
  height=#1,
  breakable
}

% Hint box
\newtcolorbox{hintbox}{
  colback=solutiongreen!5,
  colframe=solutiongreen,
  fonttitle=\bfseries\color{white},
  title=Hinweis,
  breakable
}

% Julia code style
\lstdefinelanguage{Julia}{
  keywords={function, end, return, if, else, elseif, for, while, in, begin, let, do, try, catch, finally, using, import, export, mutable, struct, abstract, type},
  keywordstyle=\color{theoryblue}\bfseries,
  comment=[l]{\#},
  commentstyle=\color{hintgray}\itshape,
  stringstyle=\color{solutiongreen},
  morestring=[b]",
  sensitive=true
}
\lstset{
  language=Julia,
  backgroundcolor=\color{codebg},
  basicstyle=\ttfamily\small,
  breaklines=true,
  frame=single,
  rulecolor=\color{black!30},
  numbers=left,
  numberstyle=\tiny\color{hintgray}
}

\pagestyle{fancy}
\fancyhf{}
\rhead{\textcolor{hintgray}{\small Rocket Science Workbook}}
\lhead{\textcolor{hintgray}{\small \leftmark}}
\cfoot{\thepage}

\titleformat{\section}{\large\bfseries\color{theoryblue}}{}{0em}{}[\titlerule]
\titleformat{\subsection}{\normalsize\bfseries\color{exerciseorange}}{}{0em}{}


\begin{document}

\begin{center}
  {\LARGE \textbf{Kapitel 05: Der Gravity Turn}}\\[0.5em]
  {\large \textcolor{hintgray}{Wie Raketen ihren Pitchwinkel steuern und in einen Orbit gelangen}}\\[0.3em]
  {\small \textcolor{hintgray}{Rocket Science Workbook — simulation-2D-jl}}
\end{center}
\vspace{1em}
\hrule
\vspace{1.5em}

\section{Was ist ein Gravity Turn?}

\begin{theorybox}{Grundprinzip}
Beim \textbf{Gravity Turn} wird die Rakete nicht aktiv um eine Achse gesteuert. Stattdessen:
\begin{enumerate}
  \item Die Rakete startet \textbf{senkrecht} ($\theta = 90°$)
  \item Nach einer kurzen Vertikalphase wird ein kleiner \textbf{Kick-Winkel} eingeleitet
  \item Die Schwerkraft dreht den Geschwindigkeitsvektor langsam zur Horizontale
  \item Der Schubvektor folgt dem Geschwindigkeitsvektor (minimaler Anstellwinkel)
\end{enumerate}
\textbf{Vorteil:} Minimaler aerodynamischer Stress. Die Rakete ``fällt'' in ihre Trajektorie.
\end{theorybox}

\section{Pitchwinkel-Strategien}

\begin{theorybox}{Strategie 1: Lineares Pitchprogramm}
\[
\theta(t) =
\begin{cases}
  90° & t \leq t_\text{kick} \\
  90° - \dfrac{90° - \theta_\text{end}}{t_\text{end} - t_\text{kick}}\,(t - t_\text{kick}) & t_\text{kick} < t \leq t_\text{end} \\
  \theta_\text{end} & t > t_\text{end}
\end{cases}
\]
Einfach, aber nicht optimal. Gut zum Lernen.
\end{theorybox}

\begin{theorybox}{Strategie 2: Velocity-Vector Tracking}
Der Pitchwinkel folgt dem \textbf{Geschwindigkeitsvektor}:
\[
  \theta(t) = \arctan\!\left(\frac{v_y}{v_x}\right)
\]
Das ist der ``echte'' Gravity Turn — die Rakete richtet sich immer in Flugrichtung aus.
\end{theorybox}

\section{Die Tsiolkovsky-Raketengleichung}

\begin{theorybox}{Ideales $\Delta v$}
Ohne Schwerkraft und Drag gilt:
\[
  \Delta v = v_e \ln\!\left(\frac{m_0}{m_f}\right)
\]
$v_e$: effektive Ausströmgeschwindigkeit, $m_0$: Startmasse, $m_f$: Endmasse (nach Treibstoffverbrauch).\\
Das ist eine \textbf{obere Schranke} für das erreichbare $\Delta v$ — Drag und Gravitationsverluste reduzieren es.
\end{theorybox}

\section{Rechenaufgaben}

\begin{exercisebox}{1 — Tsiolkovsky-Gleichung}
Gegeben: $m_0 = 100\,000\,\text{kg}$, $m_f = 10\,000\,\text{kg}$, $v_e = 3200\,\text{m/s}$.\\[0.5em]
\textbf{a)} Berechne das ideale $\Delta v$:
\[
  \Delta v = 3200 \cdot \ln\!\left(\frac{100\,000}{10\,000}\right) = \underline{\hspace{4cm}}\,\text{m/s}
\]
\textbf{b)} Die erste kosmische Geschwindigkeit (LEO) beträgt $\approx 7\,800\,\text{m/s}$.\\
Reicht eine Stufe? Was hemmt uns in der Praxis?\\[0.5em]
\textbf{c)} Mit zwei gleichen Stufen (je $m_0 = 50\,000$, $m_f = 5\,000$, $v_e = 3200$):
\[
  \Delta v_\text{total} = \Delta v_1 + \Delta v_2 = \underline{\hspace{6cm}}\,\text{m/s}
\]
\end{exercisebox}

\begin{solutionspace}[7cm]
\end{solutionspace}

\begin{exercisebox}{2 — Lineares Pitchprogramm berechnen}
Parameter: $t_\text{kick} = 15\,\text{s}$, $\theta_\text{end} = 0°$, $t_\text{end} = 150\,\text{s}$.\\[0.5em]
Berechne $\theta(t)$ für:

\begin{center}
\begin{tabular}{|c|c|}
\hline
$t$ (s) & $\theta$ (°) \\
\hline
0   &  \\
\hline
15  &  \\
\hline
50  &  \\
\hline
80  &  \\
\hline
120 &  \\
\hline
150 &  \\
\hline
\end{tabular}
\end{center}

\textbf{Zeichne} die Kurve $\theta(t)$ skizzenhaft in das folgende Koordinatensystem:

\begin{center}
\begin{tikzpicture}
\begin{axis}[
  width=10cm, height=5cm,
  xlabel={$t$ (s)},
  ylabel={$\theta$ (°)},
  xmin=0, xmax=160,
  ymin=-5, ymax=100,
  grid=both,
  grid style={line width=0.3pt, draw=gray!30},
  major grid style={line width=0.6pt, draw=gray!60}
]
\end{axis}
\end{tikzpicture}
\end{center}
\end{exercisebox}

\begin{solutionspace}[3cm]
\end{solutionspace}

\begin{exercisebox}{3 — Velocity-Vector Tracking}
Eine Rakete hat $v_x = 800\,\text{m/s}$, $v_y = 200\,\text{m/s}$.\\[0.5em]
\textbf{a)} Berechne den Pitchwinkel beim Velocity-Vector Tracking:
\[
  \theta = \arctan\!\left(\frac{v_y}{v_x}\right) = \arctan\!\left(\frac{200}{800}\right) = \underline{\hspace{3cm}}°
\]
\textbf{b)} Was bedeutet $\theta \approx 14°$? Beschreibe die Flugrichtung qualitativ.\\[0.5em]
\textbf{c)} Bei welchem Verhältnis $v_y / v_x$ ist $\theta = 45°$?
\end{exercisebox}

\begin{solutionspace}[5cm]
\end{solutionspace}

\begin{exercisebox}{4 — Gravitationsverlust}
Der ``Gravitationsverlust'' (gravity drag) ist:
\[
  \Delta v_\text{grav} = \int_0^{t_\text{MECO}} g(t)\,\sin\theta(t)\,dt
\]
\textbf{Näherung} für konstantes $g = 9{,}81$ und lineares Pitch von $90°$ auf $0°$ in $t = 150\,\text{s}$:\\[0.5em]
\[
  \Delta v_\text{grav} \approx g \int_0^{150} \sin\!\left(90° - \frac{90°}{150}\,t\right)\,dt
\]
\textbf{a)} Vereinfache das Integral analytisch (Substitution oder direkt).\\[0.5em]
\textbf{b)} Berechne das Ergebnis numerisch (Trapezregel mit 3 Stützstellen: $t = 0, 75, 150$).
\end{exercisebox}

\begin{solutionspace}[8cm]
\end{solutionspace}

\section{Verbindung zum Julia-Template}

\begin{theorybox}{Was du jetzt implementierst}
In \texttt{template.jl} (Kapitel 05) implementierst du:
\begin{enumerate}
  \item \texttt{pitch\_linear(t, params)} — lineares Pitchprogramm
  \item \texttt{pitch\_velocity\_tracking(t, s)} — Velocity-Vector Tracking
  \item Vollständige 2D-Simulation mit beiden Pitch-Strategien
  \item \textbf{Trajektorie-Plot}: $x$ vs.\ $y$ (Downrange vs.\ Höhe)
  \item Vergleich: Wie viel Downrange-Distanz bei gleicher Treibstoffmenge?
\end{enumerate}
\end{theorybox}

\end{document}
